% !TEX root = ../thesis.tex

\chapter{Evaluation Framework}
\label{ch:evaluation}

% Expansion of proposal §5.3 (Evaluation Metrics). Methodology only — every
% definition, formula, and rationale lives here. Numerical results live in
% Chapter~\ref{ch:results}.

\section{Two Pillars: Utility and Privacy}
\label{sec:eval:pillars}

% TODO: One short paragraph framing the two-pillar structure adopted in this
% thesis, anchored in \cite{hernandez2022synthetic,Murtaza2023}.
%  - Utility: encompasses both statistical similarity (do the synthetic
%    distributions resemble the real ones?) and predictive performance
%    (do downstream models trained on synthetic perform comparably?).
%    These two branches are complementary: similarity is necessary but not
%    sufficient for downstream task success, and high predictive
%    performance without distributional match can mask shortcut behaviour.
%  - Privacy: does the synthetic data leak information about real records?

This chapter defines a two-pillar evaluation framework: \emph{utility} and
\emph{privacy}. The utility pillar is treated as a single concept with two
complementary branches: \emph{statistical similarity} between the synthetic
and real datasets (Section~\ref{sec:eval:utility:similarity}) and
\emph{predictive performance} of downstream models trained on the synthetic
data (Section~\ref{sec:eval:utility:predictive}). The privacy pillar is
defined in Section~\ref{sec:eval:privacy} via instance-level distance and
inference-attack metrics. Section~\ref{sec:eval:aggregate} introduces the
aggregation views that combine both pillars into a single comparative
picture.

\section{Utility Metrics}
\label{sec:eval:utility}

% Maps to data_analysis.ipynb cells 29–41.

The utility pillar combines two complementary aspects. The first is whether
the synthetic dataset \emph{looks like} the real one in a distributional
sense; the second is whether the synthetic dataset \emph{acts like} the real
one when used as training data for downstream predictive and survival models.
The same set of synthetic records is evaluated under both views.

\subsection{Statistical Similarity to the Real Data}
\label{sec:eval:utility:similarity}

% Maps to data_analysis.ipynb cells 29–34. These metrics correspond to what
% \cite{Murtaza2023,hernandez2022synthetic} term resemblance / fidelity; in
% this thesis they are the first branch of the utility pillar.

\subsubsection{Per-Column Distributional Match}
\label{sec:eval:utility:similarity:dist}

% TODO: KDE for continuous, bar plots for categorical. Visual diagnostic only.

\subsubsection{Summary Statistics (Mean, Standard Deviation)}
\label{sec:eval:utility:similarity:stats}

% TODO: Per-column mean / std table; relative differences from real.

\subsubsection{Pairwise Correlation Difference (PCD)}
\label{sec:eval:utility:similarity:pcd}

% Following \cite{goncalves2020generation}, define
%   PCD = || corr(Real) - corr(Synth) ||_F
% the Frobenius norm of the Pearson correlation matrix difference.
% Smaller PCD indicates better preservation of joint linear structure.

\subsubsection{Kolmogorov--Smirnov Two-Sample Test}
\label{sec:eval:utility:similarity:ks}

% TODO:
%  - Definition of KS statistic.
%  - Per-column KS computed via scipy.stats.ks_2samp.
%  - Note: the proposal listed KL divergence; the implemented pipeline uses
%    KS as a non-parametric alternative requiring no PMF binning. Discuss
%    in §7.5.5 (Limitations) and Ch 2 (background) — both metrics share the
%    same purpose of distributional similarity.

\subsubsection{PCA-Based Manifold Overlap}
\label{sec:eval:utility:similarity:pca}

% Standardise (StandardScaler), fit PCA on combined Real + Synth, project to
% 2D, scatter plot per source. Visual assessment of manifold overlap.

\subsection{Predictive Performance of Downstream Models}
\label{sec:eval:utility:predictive}

% Maps to data_analysis.ipynb cells 35–41.

\subsubsection{Train--Test Regimes: TRTR and TSTR}
\label{sec:eval:utility:predictive:tstr}

% TRTR (train on real, test on real) baseline; TSTR (train on synthetic,
% test on real) per generator. The proposal additionally listed TRTS and
% TSTS; these are scope-narrowed in the implemented pipeline. Note this in
% §7.5.5 (Limitations).

\subsubsection{Classifier Suite}
\label{sec:eval:utility:predictive:classifiers}

% Three downstream classifiers:
%  - Logistic Regression (max_iter=1000)
%  - Random Forest (n_estimators=100)
%  - Gradient Boosting (n_estimators=100)
% Justification: covers a linear baseline, a nonparametric ensemble, and a
% boosted tree learner — capturing diverse decision-boundary biases.

\subsubsection{Performance Metrics}
\label{sec:eval:utility:predictive:metrics}

% Per (Train Source × Classifier) cell, report:
%  - Accuracy
%  - F1 (zero_division=0)
%  - AUC-ROC
%  - Precision (zero_division=0)
%  - Recall (zero_division=0)
% AUC-ROC is the headline metric (class-imbalance robust); F1 secondary.

\subsubsection{ROC Curve Analysis}
\label{sec:eval:utility:predictive:roc}

% Gradient-Boosting ROC curves overlaid for Real (TRTR) vs. each synthetic
% source. Compare AUC and the false-positive trade-off.

\subsubsection{Feature Importance Preservation}
\label{sec:eval:utility:predictive:fi}

% Random Forest .feature_importances_ ranking per training source.
% Comparison metric: Spearman rank correlation of feature importances
% (Real vs. each Synth). Higher ρ = better preservation.

\subsubsection{Cox Proportional Hazards Regression}
\label{sec:eval:utility:predictive:cox}

% Cox model from lifelines.
%  - Duration: time_in_hospital
%  - Event:    readmitted
%  - Covariates: age, gender, num_medications, num_lab_procedures,
%    num_procedures, number_diagnoses, number_inpatient, number_emergency,
%    number_outpatient, insulin, diabetesMed, change, admission_type_id.
% Reports hazard ratios with 95 % confidence intervals.
% Note: this analysis is an addition beyond the proposal — it provides
% medically interpretable utility evidence.

\subsubsection{Cox CI Overlap Analysis}
\label{sec:eval:utility:predictive:cox-overlap}

% For each covariate, compute fraction of overlap between Real-fit CI and
% Synth-fit CI. Aggregate fraction (covariates with overlapping CI) is a
% scalar utility score: higher = better preservation of causal/risk
% relationships.

\section{Privacy Metrics}
\label{sec:eval:privacy}

% Maps to data_analysis.ipynb cells 42–48.

\subsection{Nearest-Neighbour Distance}
\label{sec:eval:privacy:nn}

% For each synthetic record, compute the L2 distance (in StandardScaler
% feature space) to its nearest real record. Report median NN distance per
% generator. Larger = better privacy (no near-duplicates of real patients).
% Real→Real baseline: second-nearest-neighbour distance within real data
% (first nearest is self-match).

\subsection{Nearest-Neighbour Distance Ratio (NNDR)}
\label{sec:eval:privacy:nndr}

% Define
%   NNDR = median(d_{Synth → Real}) / median(d_{Real → Real})
% NNDR > 1 indicates synthetic records are farther from real than real
% records are from each other — the desired privacy regime.

\subsection{Nearest-Neighbour Adversarial Accuracy (NNAA)}
\label{sec:eval:privacy:nnaa}

% Yale et al. \cite{healthgan} adversarial-accuracy test:
%   AA = 0.5 × (mean(d_{TS} > d_{TT}) + mean(d_{ST} > d_{SS}))
% computed separately for the train and test halves of real data.
% Privacy loss = TestAA - TrainAA (target ≈ 0).
% Promote NNAA from a literature-review citation to a first-class metric.

\subsection{Membership Inference Attack}
\label{sec:eval:privacy:mia}

% Threat model: adversary has access to synthetic data and a candidate
% record; must decide whether the record was in the training set.
%
% Adversary: distance-feature logistic regression
%  Features per record: min, mean, max distance to synthetic records
%  Labels: 1 = in training set, 0 = in held-out test set
%  Sample size: min(5,000 train members, 5,000 test non-members).
%  Evaluation: 5-fold cross-validated ROC-AUC.
%
% Reported metrics: MIA AUC (target ≈ 0.5, random) and MIA Advantage
% (= AUC - 0.5; target ≈ 0).

\subsection{Attribute Inference Attack}
\label{sec:eval:privacy:aia}

% Threat model: adversary knows all but one sensitive attribute of a target
% record and tries to infer the missing attribute by querying the synthetic
% data.
%
% Method: k-NN lookup in synthetic data (k = 10), majority vote among
% neighbours' values for the target attribute.
%
% Sensitive attributes evaluated: readmitted, gender, age, diag_1.
% Baseline: random-guess accuracy (= class prior).
% Advantage = Accuracy_attack - Accuracy_baseline.
% Reported metric: max attribute-inference advantage across the four
% attributes (target ≈ 0).

\subsection{Clustering Proximity / Outlier Linkability}
\label{sec:eval:privacy:linkability}

% Identify outliers in real data: top 5 % of records ranked by distance from
% the centroid in StandardScaler feature space. Compute the fraction of
% synthetic records within a privacy radius of any outlier, where the radius
% is the median real→real NN distance. Lower = better privacy.
% Note: the proposal §5.3.2 (item 4 + item 5) cited a DBSCAN-based outlier
% definition \cite{wang2019generating}; the implemented pipeline uses the
% centroid-distance variant for compute simplicity. Discuss in §7.5.5.

\section{Comparative Aggregation}
\label{sec:eval:aggregate}

% Maps to data_analysis.ipynb cells 49–52.

\subsection{Privacy--Utility Tradeoff Scatter}
\label{sec:eval:aggregate:tradeoff}

% Single 2D plot:
%   X = utility (Gradient-Boosting AUC-ROC under TSTR)
%   Y = privacy (NNDR median; higher = better)
% Reference line at NNDR = 1.0 (privacy threshold) and at TRTR AUC.

\subsection{Six-Dimension Radar Chart}
\label{sec:eval:aggregate:radar}

% Six axes (all scaled to [0, 1] with higher = better). Three axes belong to
% the utility pillar (similarity + predictive); three to the privacy pillar.
%
%  Utility axes:
%   1. Statistical similarity:    1 - mean(|corr difference|)
%   2. Predictive performance:    Gradient-Boosting AUC-ROC
%   3. Cox CI overlap:            fraction of covariates with overlapping CIs
%
%  Privacy axes:
%   4. Privacy (NN):              NNDR median, capped to [0, 1] for plotting
%   5. Privacy (MIA):             1 - |MIA AUC - 0.5|
%   6. Privacy (AIA):             1 - max(attr inference advantage)

\subsection{Comprehensive Summary Table}
\label{sec:eval:aggregate:table}

% 18 metrics × 5 sources (Real (TRTR) + 4 generators). Columns are grouped
% by pillar:
%
%  Utility / Statistical similarity:    mean KS, mean |corr diff|
%  Utility / Predictive performance:    AUC, F1, Cox CI overlap %
%  Privacy / Distance:                  NN dist median, NNDR median
%  Privacy / Membership inference:      MIA AUC, MIA Advantage
%  Privacy / Attribute inference:       max attr inf advantage
%  Privacy / Linkability:               outlier linkability %
%
% Conventions: ↑ better / ↓ better / ≈ 0.5 best — denote in the column header.

\cleardoublepage
